\subsection{Ablation Study}\label{sec:ablation}
We conduct an ablation study over multiple variants on the same FC-Siam-diff-style backbone: baseline, attention-only, \cosa{} with multi-scale only, \cosa{} with gate only, \cosa{} with single-scale and fixed gate, alignment-first, and the full \cosa{} configuration with multi-scale refinement and learnable gating.

Table~\ref{tab:ablation} lists quantitative results for each variant on LEVIR-CD. Figure~\ref{fig:ablation_qualitative} complements Table~\ref{tab:ablation} with qualitative error maps (false positives and false negatives) for the same setting.

% Table 2: CoSA ablation study (variants / components)
% Source: ablation/ablation_study/ablation_summary.json
\begin{table*}[t]
  \centering
  \caption{Ablation study on LEVIR-CD (changed class).}
  \label{tab:ablation}
  \setlength{\tabcolsep}{6pt}
  \renewcommand{\arraystretch}{1.05}
  \small
  \begin{tabularx}{\textwidth}{X r r r r}
    \toprule
    Variant & \PrecC{} & \RecC{} & \Fonec{} & \Iouc{} \\
    \midrule
    Baseline & 90.23 & 85.31 & 87.70 & 78.09 \\
    Attention-only & 90.27 & 85.61 & 87.88 & 78.38 \\
    \cosa{} with multi-scale only & 91.35 & 83.68 & 87.35 & 77.54 \\
    \cosa{} with gate only & 93.33 & 82.67 & 87.68 & 78.06 \\
    \cosa{} with single-scale and fixed gate & \textbf{93.88} & 83.22 & 88.23 & 78.94 \\
    Alignment-first & 92.00 & 85.62 & 88.70 & 79.69 \\
    Full \cosa{} (multi-scale + learnable gate; ours) & 90.39 & \textbf{88.52} & \textbf{89.45} & \textbf{80.91} \\
    \bottomrule
  \end{tabularx}
\end{table*}


\begin{figure}[!t]
  \centering
  \includegraphics[width=\columnwidth]{ablation_qualitative_alt3_errors_only_v2.png}
  \caption{Qualitative ablation comparison (errors only): highlighted regions emphasize false positives and false negatives across variants.}
  \label{fig:ablation_qualitative}
\end{figure}

``Multi-scale'' inserts \cosa{} at two resolutions (1/8 and 1/16). ``Learnable gate'' uses a trainable residual scale $\lambda$ (Section~\ref{sec:method}) initialized near zero; fixed-scale uses a constant residual factor. ``Gate only'' disables multi-scale placement and keeps only residual gating at the bottleneck. ``Alignment-first'' estimates local displacement and warps features before differencing. The full \cosa{} configuration (multi-scale + learnable residual scale) gives the best F1/IoU balance among tested variants. Relative to baseline, \cosa{} improves recall (+3.21 points) with near-stable precision (+0.16 points), yielding +1.75 \Fonec{} and +2.82 \Iouc{}. The single-scale fixed variant improves precision but reduces recall, indicating that correlation refinement without adaptive residual control can become overly conservative. The alignment-first variant improves over baseline but remains below full \cosa{}, suggesting that adaptive correlation-guided refinement is more effective than geometric pre-alignment alone in this setup \cite{shanshan2024,yan2024}. Table~\ref{tab:ablation} therefore provides the experimental evidence for the method-level claim: multi-scale sampling alone or gating alone is insufficient for peak performance, and the full gain appears only when correlation-guided sampling, multi-scale fusion, and adaptive residual control are used together. This combination enables a form of targeted refinement that is not achievable by conventional attention or single-component designs. Figure~\ref{fig:ablation_qualitative} illustrates where partial variants fail: the full module suppresses spurious detections while recovering missing change regions more consistently than partial variants.
